\documentclass[aspectratio=169]{beamer}

% - Кодировка и язык -
\usepackage[T2A]{fontenc}
\usepackage[utf8]{inputenc}
\usepackage[russian]{babel}

% - Формулы -
\usepackage{amsmath,amssymb}

% - Тема -
\usetheme{Madrid}
\setbeamertemplate{navigation symbols}{}
\setbeamertemplate{bibliography item}[text]

% - Титульная информация -
\title[Линеаризация внимания в ViT-сегментаторах]{Линеаризация механизма внимания в ViT-сегментаторах для построения динамических зон интереса}
\author[Салимли, Лобанов, Черепнов]{Салимли~А.~М., Лобанов~П.~М., Черепнов~М.~М.}
\institute[СПбПУ]{
  Санкт-Петербургский политехнический университет Петра Великого\\[2pt]
  Математика и компьютерные науки\\
  Искусственный интеллект и машинное обучение\\[2pt]
  гр.~5140201/60301
}
\date{2026}

\begin{document}

% - Титульный слайд -
\begin{frame}
  \titlepage
\end{frame}

% - Введение -
\begin{frame}{Введение}
  \small
  \begin{itemize}
    \item CNN: ограниченное рецептивное поле, нет глобального контекста, сегментация сводится к классификации пикселей
    \item ViT: глобальный контекст за счёт самовнимания, но признаки одномасштабные
    \item Решение ?: SegFormer - многоуровневые признаки, Efficient Self-Attention снижает сложность до $O(N^2/R)$, но она остаётся квадратичной
    \item Научная проблема - точность глобального внимания против его квадратичной сложности, мешающей работе в реальном времени
    \item Объект - сегментаторы на основе ViT; предмет - методы линеаризации механизма внимания
    \item Цель - найти метод линеаризации, сохраняющий точность сегментации и обеспечивающий работу в реальном времени
    \item Задачи - анализ методов линеаризации, сравнение по точности и быстродействию, адаптация к многоуровневой архитектуре, экспериментальная апробация
    \item Применение - КТ и другие медицинские снимки, дорожные сцены, контроль дефектов, динамическая зона интереса трамвайных путей~\cite{myarticle}
  \end{itemize}
\end{frame}

% - Список источников -
\begin{frame}{Список источников}
  \small
  \begin{thebibliography}{9}
    \bibitem{myarticle} Салимли~А.~М., Востров~А.~В. Построение зоны интереса на основе сегментации // Неделя науки ИКНК. - 2026. - С.~386.
  \end{thebibliography}
\end{frame}

\end{document}
